\chapter{2020年天津卷}

\section{单选题}

\begin{problem}
设全集 \(U=\{-3,-2,-1,0,1,2,3\}\)，集合 \(A=\{-1,0,1,2\}\)，\(B=\{-3,0,2,3\}\)，则 \(A\cap\left( U\setminus B \right)=\)（\quad）

\choices
    {\(\{-3,3\}\)}
    {\(\{0,2\}\)}
    {\(\{-1,1\}\)}
    {\(\{-3,-2,-1,1,3\}\)}

\end{problem}

\begin{answer}
C
\end{answer}

\begin{solution}
由题意结合补集的定义可知： \( U\setminus B=\{-2,-1,1\} \)，则 \( A\cap\left( U\setminus B \right)=\{-1,1\} \) .

故选：C.
\end{solution}
\begin{problem}
设 \(a\in\R\)，则“\(a>1\)”是“\(a^2>a\)”的（\quad）

\choices
    {充分不必要条件}
    {必要不充分条件}
    {充要条件}
    {既不充分也不必要条件}

\end{problem}

\begin{answer}
A
\end{answer}

\begin{solution}
求解二次不等式 \( a^2>a \) 可得： \( a>1 \) 或 \( a<0 \) ，

据此可知： \( a>1 \) 是 \( a^2>a \) 的充分不必要条件.

故选：A.
\end{solution}
\begin{problem}
函数 \(y=\dfrac{4x}{x^2+1}\) 的图象大致为（\quad）

\choices
    {\bitmapinclude[max width=3.3cm,max totalheight=4.8cm]{img/2020/tianjin/q3_fig1.png}}
    {\bitmapinclude[max width=3.3cm,max totalheight=4.8cm]{img/2020/tianjin/q3_fig2.png}}
    {\bitmapinclude[max width=3.3cm,max totalheight=4.8cm]{img/2020/tianjin/q3_fig3.png}}
    {\bitmapinclude[max width=3.3cm,max totalheight=4.8cm]{img/2020/tianjin/q3_fig4.png}}

\end{problem}

\begin{answer}
A
\end{answer}

\begin{solution}
由函数的解析式可得： \( f(-x)=\dfrac{-4x}{x^2+1}=-f(x) \) ，则函数 \( f(x) \) 为奇函数，其图象关于坐标原点对称，选项 \( C \) ，\( D \) 错误；

当 \( x=1 \) 时， \( y=\dfrac{4}{1+1}=2>0 \) ，选项 \( B \) 错误.

故选：A.
\end{solution}
\begin{problem}
从一批零件中抽取 80 个，测量其直径（单位：mm），将所得数据分为 9 组 \([5.31,5.33),[5.33,5.35),\cdots,[5.45,5.47),[5.47,5.49]\)，并整理得到如下频率分布直方图，则在被抽取的零件中，直径落在区间 \([5.43,5.47)\) 内的个数为（\quad）

\begin{center}
\bitmapfigure[max width=0.90\linewidth,max totalheight=4.8cm]{img/2020/tianjin/q4_fig1.png}
\end{center}

\choices
    {10}
    {18}
    {20}
    {36}

\end{problem}

\begin{answer}
B
\end{answer}

\begin{solution}
根据直方图，直径落在区间 \([5.43,5.47)\) 之间的零件频率为： \( \left( 6.25+5.00 \right)\times0.02=0.225 \) ，

则区间 \([5.43,5.47)\) 内零件的个数为： \( 80\times0.225=18 \) .

故选：B.
\end{solution}
\begin{problem}
若棱长为 \(2\sqrt3\) 的正方体的顶点都在同一球面上，则该球的表面积为（\quad）

\choices
    {\(12\pi\)}
    {\(24\pi\)}
    {\(36\pi\)}
    {\(144\pi\)}

\end{problem}

\begin{answer}
C
\end{answer}

\begin{solution}
这个球是正方体的外接球，其半径等于正方体的体对角线的一半，
\(R=\frac{\sqrt{\left( 2\sqrt3 \right)^2+\left( 2\sqrt3 \right)^2+\left( 2\sqrt3 \right)^2}}{2}=3\).

所以，这个球的表面积为 \( S=4\pi R^2=4\pi\times3^2=36\pi \) .

故选：C.
\end{solution}
\begin{problem}
设 \(a=3^{0.7}\)，\(b=\left( \dfrac13 \right)^{-0.8}\)，\(c=\log_{0.7}0.8\)，则 \(a,b,c\) 的大小关系为（\quad）

\choices
    {\(a<b<c\)}
    {\(b<a<c\)}
    {\(b<c<a\)}
    {\(c<a<b\)}

\end{problem}

\begin{answer}
D
\end{answer}

\begin{solution}
因为 \( a=3^{0.7}>1 \) ，

\(b=\left( \frac13 \right)^{-0.8}=3^{0.8}>3^{0.7}=a\)，

\(c=\log_{0.7}0.8<\log_{0.7}0.7=1\)，

所以 \( c<1<a<b \) .

故选：D.
\end{solution}
\begin{problem}
设双曲线 \(C\) 的方程为 \(\dfrac{x^2}{a^2}-\dfrac{y^2}{b^2}=1\)（\(a>0\)，\(b>0\)），过抛物线 \(y^2=4x\) 的焦点和点 \((0,b)\) 的直线为 \(l\). 若 \(C\) 的一条渐近线与 \(l\) 平行，另一条渐近线与 \(l\) 垂直，则双曲线 \(C\) 的方程为（\quad）

\choices
    {\(\dfrac{x^2}{4}-\dfrac{y^2}{4}=1\)}
    {\(x^2-\dfrac{y^2}{4}=1\)}
    {\(\dfrac{x^2}{4}-y^2=1\)}
    {\(x^2-y^2=1\)}

\end{problem}

\begin{answer}
D
\end{answer}

\begin{solution}
由题可知，抛物线的焦点为 \((1,0)\) ，所以直线 \( l \) 的方程为 \( x+\dfrac{y}{b}=1 \) ，即直线的斜率为 \(-b\) ，

又双曲线的渐近线方程为 \( y=\pm\dfrac{b}{a}x \) ，所以 \(-b=-\dfrac{b}{a}\) ，\(-b\times\dfrac{b}{a}=-1\) ，因为 \( a>0 \) ，\( b>0 \) ，解得 \( a=1 \) ，\( b=1 \) .

故选：D.
\end{solution}
\begin{problem}
已知函数 \(f(x)=\sin\left( x+\dfrac{\pi}{3} \right)\). 给出下列结论：

\begin{circlelist}
\item \(f(x)\) 的最小正周期为 \(2\pi\)；

\item \(f\left( \dfrac{\pi}{2} \right)\) 是 \(f(x)\) 的最大值；

\item 把函数 \(y=\sin x\) 的图象上所有点向左平移 \(\dfrac{\pi}{3}\) 个单位长度，可得到函数 \(y=f(x)\) 的图象.
\end{circlelist}

其中所有正确结论的序号是（\quad）

\choices
    {\(\circled{1}\)}
    {\(\circled{1}\circled{3}\)}
    {\(\circled{2}\circled{3}\)}
    {\(\circled{1}\circled{2}\circled{3}\)}

\end{problem}

\begin{answer}
B
\end{answer}

\begin{solution}
因为 \( f(x)=\sin\left( x+\dfrac{\pi}{3} \right) \) ，所以周期
\(T=\frac{2\pi}{\omega}=2\pi\)，
故 \circled{1} 正确；
因为
\(f\left( \frac{\pi}{2} \right)=\sin\left( \frac{\pi}{2}+\frac{\pi}{3} \right)=\sin\frac{5\pi}{6}=\frac12\neq1\)，
故 \circled{2} 不正确；

将函数 \( y=\sin x \) 的图象上所有点向左平移 \( \dfrac{\pi}{3} \) 个单位长度，得到 \( y=\sin\left( x+\dfrac{\pi}{3} \right) \) 的图象，故 \circled{3} 正确.

故选：B.
\end{solution}
\begin{problem}
已知函数 \(f(x)=x^3\)（\(x\ge 0\)），\(f(x)=-x\)（\(x<0\)）. 若函数 \(g(x)=f(x)-\left| kx^2-2x \right|\)（\(k\in\R\)）恰有 4 个零点，则 \(k\) 的取值范围是（\quad）

\choices
    {\(\left( -\infty,-\dfrac12 \right)\cup\left( 2\sqrt2,+\infty \right)\)}
    {\(\left( -\infty,-\dfrac12 \right)\cup\left( 0,2\sqrt2 \right)\)}
    {\(\left( -\infty,0 \right)\cup\left( 0,2\sqrt2 \right)\)}
    {\(\left( -\infty,0 \right)\cup\left( 2\sqrt2,+\infty \right)\)}

\end{problem}

\begin{answer}
D
\end{answer}

\begin{solution}
注意到 \( g(0)=0 \) ，所以要使 \( g(x) \) 恰有 4 个零点，只需考察 \( g(x)=0 \) 的其余非零零点. 对于 \( x\neq0 \) ，由 \( g(x)=0 \) 可得
\(\left| kx-2 \right|=\frac{f(x)}{|x|}\)
反之，方程 \( \left| kx-2 \right|=\dfrac{f(x)}{|x|} \) 的每一个实根也对应 \( g(x)=0 \) 的一个非零零点. 因此，函数 \( g(x) \) 恰有 4 个零点，当且仅当方程 \( \left| kx-2 \right|=\dfrac{f(x)}{|x|} \) 恰有 3 个实根.

令
\(h(x)=\frac{f(x)}{|x|}\)，
即 \( y=\left| kx-2 \right| \) 与 \( y=h(x) \) 的图象有 3 个不同交点.

因为
\(h(x)=x^2\text{（}x>0\text{）}\)，\(h(x)=1\text{（}x<0\text{）}\).

当 \( k=0 \) 时，此时 \( y=2 \) ，如图 1，\( y=2 \) 与 \( y=h(x) \) 有 2 个不同交点，不满足题意；

\begin{center}
\bitmapfigure[max width=0.50\linewidth,max totalheight=4.8cm]{img/2020/tianjin/q9_solution_fig1.png}
\end{center}

当 \( k<0 \) 时，如图 2，此时 \( y=\left| kx-2 \right| \) 与 \( y=h(x) \) 恒有 3 个不同交点，满足题意；

\begin{center}
\bitmapfigure[max width=0.50\linewidth,max totalheight=4.8cm]{img/2020/tianjin/q9_solution_fig3.png}
\end{center}

当 \( k>0 \) 时，如图 3，当 \( y=kx-2 \) 与 \( y=x^2 \) 相切时，联立方程得
\(x^2-kx+2=0\)，
令 \( \Delta=0 \) 得 \( k^2-8=0 \) ，解得 \( k=2\sqrt2 \)（负值舍去），所以 \( k>2\sqrt2 \) .

\begin{center}
\bitmapfigure[max width=0.50\linewidth,max totalheight=4.8cm]{img/2020/tianjin/q9_solution_fig2.png}
\end{center}

综上， \( k \) 的取值范围为 \( \left( -\infty,0 \right)\cup\left( 2\sqrt2,+\infty \right) \) .

故选：D.
\end{solution}
\section{填空题}

\begin{problem}
\(\i\) 是虚数单位，复数 \(\dfrac{8-\i}{2+\i}=\fillinblank{}\).

\end{problem}

\begin{answer}
\(3-2\i\)
\end{answer}

\begin{solution}
由有理化得
\(\frac{8-\i}{2+\i}=\frac{\left( 8-\i \right)\left( 2-\i \right)}{\left( 2+\i \right)\left( 2-\i \right)}=\frac{15-10\i}{5}=3-2\i\).

故答案为： \( 3-2\i \) .
\end{solution}
\begin{problem}
在 \(\left( x+\dfrac{2}{x^2} \right)^5\) 的展开式中，\(x^2\) 项的系数是\fillinblank{}.

\end{problem}

\begin{answer}
10
\end{answer}

\begin{solution}
因为 \( \left( x+\dfrac{2}{x^2} \right)^5 \) 的展开式的通项公式为
\(T_{r+1}=\mathrm{C}_5^r x^{5-r}\left( \frac{2}{x^2} \right)^r=\mathrm{C}_5^r\cdot2^r\cdot x^{5-3r}\)，其中
\( r=0,1,2,3,4,5 \)，
令 \( 5-3r=2 \) ，解得 \( r=1 \) .

所以 \( x^2 \) 的系数为 \( \mathrm{C}_5^1\times2=10 \) .

故答案为：10.
\end{solution}
\begin{problem}
已知直线 \(x-\sqrt3 y+8=0\) 和圆 \(x^2+y^2=r^2\)（\(r>0\)）相交于 \(A,B\) 两点. 若 \(|AB|=6\)，则 \(r=\fillinblank{}\).

\end{problem}

\begin{answer}
5
\end{answer}

\begin{solution}
因为圆心 \( (0,0) \) 到直线 \( x-\sqrt3 y+8=0 \) 的距离
\(d=\frac{8}{\sqrt{1+\left( \sqrt3 \right)^2}}=4\)，
由 \( |AB|=2\sqrt{r^2-d^2} \) 可得
\(6=2\sqrt{r^2-4^2}\)，
解得 \( r=5 \) .

故答案为：5.
\end{solution}
\begin{problem}
已知甲、乙两球落入盒子的概率分别为 \(\dfrac12\) 和 \(\dfrac13\). 假定两球是否落入盒子互不影响，则甲、乙两球都落入盒子的概率为\fillinblank{}；甲、乙两球至少有一个落入盒子的概率为\fillinblank{}.

\end{problem}

\begin{answer}
\(\dfrac16,\ \dfrac23\)
\end{answer}

\begin{solution}
甲、乙两球落入盒子的概率分别为 \( \dfrac12 \) ，\( \dfrac13 \) ，且两球是否落入盒子互不影响，

所以甲、乙都落入盒子的概率为
\(\frac12\times\frac13=\frac16\).

甲、乙两球都不落入盒子的概率为
\(\left( 1-\frac12 \right)\left( 1-\frac13 \right)=\frac13\)，
所以甲、乙两球至少有一个落入盒子的概率为
\(1-\frac13=\frac23\).

故答案为： \( \dfrac16 \) ； \( \dfrac23 \) .
\end{solution}
\begin{problem}
已知 \(a>0\)，\(b>0\)，且 \(ab=1\)，则 \(\dfrac1{2a}+\dfrac1{2b}+\dfrac8{a+b}\) 的最小值为\fillinblank{}.

\end{problem}

\begin{answer}
4
\end{answer}

\begin{solution}
因为 \( a>0 \) ，\( b>0 \) ，所以 \( a+b>0 \) ，又 \( ab=1 \) ，所以
\(\frac1{2a}+\frac1{2b}+\frac8{a+b}
=\frac{ab}{2a}+\frac{ab}{2b}+\frac8{a+b}
=\frac{a+b}{2}+\frac8{a+b}\).

由基本不等式可得
\(\frac{a+b}{2}+\frac8{a+b}\ge2\sqrt{\frac{a+b}{2}\cdot\frac8{a+b}}=4\)，
当且仅当 \( a+b=4 \) 时取等号.

结合 \( ab=1 \) ，解得 \( a=2-\sqrt3 \) ，\( b=2+\sqrt3 \) ，或 \( a=2+\sqrt3 \) ，\( b=2-\sqrt3 \) 时，等号成立.

故答案为：4.
\end{solution}
\begin{problem}
如图，在四边形 \(ABCD\) 中，\(\angle B=60^\circ\)，\(|AB|=3\)，\(|BC|=6\)，且 \(\overrightarrow{AD}=\lambda\overrightarrow{BC}\)，\(\overrightarrow{AD}\cdot\overrightarrow{AB}=-\dfrac32\)，则实数 \(\lambda\) 的值为\fillinblank{}，若 \(M,N\) 是线段 \(BC\) 上的动点，且 \(|MN|=1\)，则 \(\overrightarrow{DM}\cdot\overrightarrow{DN}\) 的最小值为\fillinblank{}.

\begin{center}
\bitmapfigure[max width=0.76\linewidth,max totalheight=4.8cm]{img/2020/tianjin/q15_fig1.png}
\end{center}

\end{problem}

\begin{answer}
\(\dfrac16,\ \dfrac{13}{2}\)
\end{answer}

\begin{solution}
因为 \( \overrightarrow{AD}=\lambda\overrightarrow{BC} \) ，所以 \( \overrightarrow{AD}//\overrightarrow{BC} \) ，故 \( \angle BAD=180^\circ-\angle B=120^\circ \) .

\(\overrightarrow{AB}\cdot\overrightarrow{AD}
=\lambda\overrightarrow{BC}\cdot\overrightarrow{AB}
=\lambda|BC||AB|\cos120^\circ
=\lambda\times6\times3\times\left( -\frac12 \right)
=-9\lambda
=-\frac32\)，
解得 \( \lambda=\dfrac16 \) .

以点 \( B \) 为坐标原点，\( BC \) 所在直线为 \( x \) 轴建立如下图所示的平面直角坐标系 \( xBy \) ，

\begin{center}
\bitmapfigure[max width=0.66\linewidth,max totalheight=4.8cm]{img/2020/tianjin/q15_solution_fig1.png}
\end{center}

则 \( C(6,0) \) .

因为 \( AB=3 \) ，\( \angle ABC=60^\circ \) ，所以
\(A\left( \frac32,\frac{3\sqrt3}{2} \right)\).

又因为 \( \overrightarrow{AD}=\dfrac16\overrightarrow{BC} \) ，所以
\(D\left( \frac52,\frac{3\sqrt3}{2} \right)\).

设 \( M(x,0) \) ，则 \( N(x+1,0) \)（其中 \( 0\le x\le5 \)），于是
\(\overrightarrow{DM}=\left( x-\frac52,-\frac{3\sqrt3}{2} \right)\)，\(\overrightarrow{DN}=\left( x-\frac32,-\frac{3\sqrt3}{2} \right)\).

所以
\(\overrightarrow{DM}\cdot\overrightarrow{DN}
=\left( x-\frac52 \right)\left( x-\frac32 \right)+\left( \frac{3\sqrt3}{2} \right)^2
=x^2-4x+\frac{21}{2}
=\left( x-2 \right)^2+\frac{13}{2}\).

所以，当 \( x=2 \) 时， \( \overrightarrow{DM}\cdot\overrightarrow{DN} \) 取得最小值 \( \dfrac{13}{2} \) .

故答案为： \( \dfrac16 \) ； \( \dfrac{13}{2} \) .
\end{solution}
\section{解答题}

\begin{problem}
在 \(\triangle ABC\) 中，角 \(A,B,C\) 所对的边分别为 \(a,b,c\). 已知 \(a=2\sqrt2\)，\(b=5\)，\(c=\sqrt{13}\).

\begin{enumerate}
\item 求角 \(C\) 的大小；
\item 求 \(\sin A\) 的值；
\item 求 \(\sin\left( 2A+\dfrac{\pi}{4} \right)\) 的值.
\end{enumerate}

\end{problem}

\begin{solution}
（1）在 \( \triangle ABC \) 中，由 \( a=2\sqrt2 \) ，\( b=5 \) ，\( c=\sqrt{13} \) 及余弦定理得
\(\cos C=\frac{a^2+b^2-c^2}{2ab} =\frac{8+25-13}{2\times2\sqrt2\times5} =\frac{\sqrt2}{2}\)，
又因为 \( C\in\left( 0,\pi \right) \) ，所以
\(C=\frac{\pi}{4}\).

（2）在 \( \triangle ABC \) 中，由 \( C=\dfrac{\pi}{4} \) ，\( a=2\sqrt2 \) ，\( c=\sqrt{13} \) 及正弦定理，可得
\(\sin A=\frac{a\sin C}{c} =\frac{2\sqrt2\times\frac{\sqrt2}{2}}{\sqrt{13}} =\frac{2\sqrt{13}}{13}\).

（3）由 \( a<c \) 知角 \( A \) 为锐角，由 \( \sin A=\dfrac{2\sqrt{13}}{13} \) ，可得
\(\cos A=\sqrt{1-\sin^2 A}=\frac{3\sqrt{13}}{13}\).

进而
\(\sin2A=2\sin A\cos A=\frac{12}{13}\)，\(\cos2A=2\cos^2 A-1=\frac{5}{13}\).

所以
\(\sin\left( 2A+\frac{\pi}{4} \right)
=\sin2A\cos\frac{\pi}{4}+\cos2A\sin\frac{\pi}{4}
=\frac{12}{13}\times\frac{\sqrt2}{2}+\frac{5}{13}\times\frac{\sqrt2}{2}
=\frac{17\sqrt2}{26}\).
\end{solution}
\begin{problem}
如图，在三棱柱 \(ABC-A_1B_1C_1\) 中，\(CC_1\perp\) 平面 \(ABC\)，\(AC\perp BC\)，\(AC=BC=2\)，\(CC_1=3\)，点 \(D,E\) 分别在棱 \(AA_1\) 和棱 \(CC_1\) 上，且 \(AD=1\)，\(CE=2\)，\(M\) 为棱 \(A_1B_1\) 的中点.



\begin{enumerate}
\item 求证：\(C_1M\perp B_1D\)；
\item 求二面角 \(B-B_1E-D\) 的正弦值；
\item 求直线 \(AB\) 与平面 \(DB_1E\) 所成角的正弦值.
\end{enumerate}

\bitmapfigure[max width=0.80\linewidth,max totalheight=4.8cm]{img/2020/tianjin/q17_fig1.png}

\end{problem}

\begin{solution}
（1）依题意，以 \( C \) 为原点，分别以 \( \overrightarrow{CA} \) ，\( \overrightarrow{CB} \) ，\( \overrightarrow{CC_1} \) 的方向为 \( x \) 轴，\( y \) 轴，\( z \) 轴的正方向建立空间直角坐标系（如图）：

\begin{center}
\bitmapfigure[max width=0.60\linewidth,max totalheight=4.8cm]{img/2020/tianjin/q17_solution_fig1.png}
\end{center}

可得
\(C(0,0,0)\)，\(A(2,0,0)\)，\(B(0,2,0)\)，\(C_1(0,0,3)\)，
\(A_1(2,0,3)\)，\(B_1(0,2,3)\)，\(D(2,0,1)\)，\(E(0,0,2)\)，\(M(1,1,3)\).

依题意，
\(\overrightarrow{C_1M}=(1,1,0)\)，\(\overrightarrow{B_1D}=(2,-2,-2)\)，
从而
\(\overrightarrow{C_1M}\cdot\overrightarrow{B_1D}=2-2+0=0\)，
所以 \( C_1M\perp B_1D \) .

（2）依题意， \( \overrightarrow{CA}=(2,0,0) \) 是平面 \( BB_1E \) 的一个法向量，

\(\overrightarrow{EB_1}=(0,2,1)\)，\(\overrightarrow{ED}=(2,0,-1)\).

	设 \( \bs{n}=(x,y,z) \) 为平面 \( DB_1E \) 的法向量，则 \( \bs{n}\cdot\overrightarrow{EB_1}=0 \)，\( \bs{n}\cdot\overrightarrow{ED}=0 \)，即 \( 2y+z=0 \)，\( 2x-z=0 \) .
	不妨设 \( x=1 \) ，可得 \( \bs{n}=(1,-1,2) \) .

于是
\(\cos\left\langle \overrightarrow{CA},\bs{n} \right\rangle
=\frac{\overrightarrow{CA}\cdot\bs{n}}{\left| \overrightarrow{CA} \right|\left| \bs{n} \right|}
=\frac{2}{2\times\sqrt6}
=\frac{\sqrt6}{6}\)，

\(\sin\left\langle \overrightarrow{CA},\bs{n} \right\rangle
=\sqrt{1-\cos^2\left\langle \overrightarrow{CA},\bs{n} \right\rangle}
=\frac{\sqrt{30}}{6}\).

所以，二面角 \( B-B_1E-D \) 的正弦值为 \( \dfrac{\sqrt{30}}{6} \) .

（3）依题意，
\(\overrightarrow{AB}=(-2,2,0)\).

由第 2 问知 \( \bs{n}=(1,-1,2) \) 为平面 \( DB_1E \) 的一个法向量，于是
\(\cos\left\langle \overrightarrow{AB},\bs{n} \right\rangle
=\frac{\overrightarrow{AB}\cdot\bs{n}}{\left| \overrightarrow{AB} \right|\left| \bs{n} \right|}
=\frac{-4}{2\sqrt2\times\sqrt6}
=-\frac{\sqrt3}{3}\).

所以，直线 \( AB \) 与平面 \( DB_1E \) 所成角的正弦值为 \( \dfrac{\sqrt3}{3} \) .
\end{solution}
\begin{problem}
已知椭圆 \(\dfrac{x^2}{a^2}+\dfrac{y^2}{b^2}=1\)（\(a>b>0\)）的一个顶点为 \(A(0,-3)\)，右焦点为 \(F\)，且 \(|OA|=|OF|\)，其中 \(O\) 为原点.

\begin{enumerate}
\item 求椭圆方程；
\item 已知点 \(C\) 满足 \(3\overrightarrow{OC}=\overrightarrow{OF}\)，点 \(B\) 在椭圆上（\(B\) 异于椭圆的顶点），直线 \(AB\) 与以 \(C\) 为圆心的圆相切于点 \(P\)，且 \(P\) 为线段 \(AB\) 的中点. 求直线 \(AB\) 的方程.
\end{enumerate}

\end{problem}

\begin{solution}
（1）因为椭圆 \( \dfrac{x^2}{a^2}+\dfrac{y^2}{b^2}=1 \)（\( a>b>0 \)）的一个顶点为 \( A(0,-3) \) ，

所以 \( b=3 \) .

由 \( OA=OF \) ，得 \( c=b=3 \) ，

又由 \( a^2=b^2+c^2 \) ，得
\(a^2=3^2+3^2=18\)，
所以，椭圆的方程为
\(\frac{x^2}{18}+\frac{y^2}{9}=1\).

（2）因为直线 \( AB \) 与以 \( C \) 为圆心的圆相切于点 \( P \) ，所以 \( CP\perp AB \) .

根据题意可知，直线 \( AB \) 和直线 \( CP \) 的斜率均存在. 设直线 \( AB \) 的斜率为 \( k \) ，则直线 \( AB \) 的方程为
\(y=kx-3\).

	由 \( y=kx-3 \) 和 \( \dfrac{x^2}{18}+\dfrac{y^2}{9}=1 \) ，消去 \( y \) ，可得
	\(\left( 2k^2+1 \right)x^2-12kx=0\)，
解得
\(x=0\)或\(x=\frac{12k}{2k^2+1}\).

将 \( x=\dfrac{12k}{2k^2+1} \) 代入 \( y=kx-3 \) ，得
\(y=\frac{6k^2-3}{2k^2+1}\)，
所以，点 \( B \) 的坐标为
\(\left( \frac{12k}{2k^2+1},\frac{6k^2-3}{2k^2+1} \right)\).

因为 \( P \) 为线段 \( AB \) 的中点，点 \( A \) 的坐标为 \( (0,-3) \) ，

所以点 \( P \) 的坐标为
\(\left( \frac{6k}{2k^2+1},\frac{-3}{2k^2+1} \right)\).

由 \( 3\overrightarrow{OC}=\overrightarrow{OF} \) ，得点 \( C \) 的坐标为 \( (1,0) \) .

所以，直线 \( CP \) 的斜率为
\(k_{CP} =\frac{\frac{-3}{2k^2+1}-0}{\frac{6k}{2k^2+1}-1} =\frac{3}{2k^2-6k+1}\).

又因为 \( CP\perp AB \) ，所以
\(k\cdot\frac{3}{2k^2-6k+1}=-1\)，
整理得
\(2k^2-3k+1=0\)，
解得 \( k=\dfrac12 \) 或 \( k=1 \) .

所以，直线 \( AB \) 的方程为
\(y=\frac12x-3\)或\(y=x-3\).
\end{solution}
\begin{problem}
已知 \(\{a_n\}\) 为等差数列，\(\{b_n\}\) 为等比数列，\(a_1=b_1=1\)，\(a_5=5\left( a_4-a_3 \right)\)，\(b_5=4\left( b_4-b_3 \right)\).

\begin{enumerate}
\item 求 \(\{a_n\}\) 和 \(\{b_n\}\) 的通项公式；
\item 记 \(\{a_n\}\) 的前 \(n\) 项和为 \(S_n\)，求证：\(S_nS_{n+2}<S_{n+1}^2\)（\(n\in\N^*\)）；
\item 对任意的正整数 \(n\)，设 \(c_n=\dfrac{(3a_n-2)b_n}{a_na_{n+2}}\)（\(n\) 为奇数），\(c_n=\dfrac{a_{n-1}}{b_{n+1}}\)（\(n\) 为偶数），求数列 \(\{c_n\}\) 的前 \(2n\) 项和.
\end{enumerate}

\end{problem}

\begin{solution}
（1）设等差数列 \( \{a_n\} \) 的公差为 \( d \) ，等比数列 \( \{b_n\} \) 的公比为 \( q \) .

由 \( a_1=1 \) ，\( a_5=5\left( a_4-a_3 \right) \) ，可得 \( d=1 \) .

从而 \( \{a_n\} \) 的通项公式为
\(a_n=n\).

由 \( b_1=1 \) ，\( b_5=4\left( b_4-b_3 \right) \) ，又 \( q\neq0 \) ，可得
\(q^4=4q^2\left( q-1 \right)\)，
即
\(q^2-4q+4=0\)，
解得 \( q=2 \) ，

从而 \( \{b_n\} \) 的通项公式为
\(b_n=2^{n-1}\).

（2）证明：由第 1 问可得
\(S_n=\frac{n(n+1)}{2}\)，

\(S_nS_{n+2}=\frac{n(n+1)(n+2)(n+3)}{4}\)，\(S_{n+1}^2=\frac{(n+1)^2(n+2)^2}{4}\).

从而
\(S_nS_{n+2}-S_{n+1}^2 =-\frac12(n+1)(n+2)<0\)，
所以
\(S_nS_{n+2}<S_{n+1}^2\).

（3）当 \( n \) 为奇数时，
\(c_n=\frac{(3a_n-2)b_n}{a_na_{n+2}}
=\frac{(3n-2)2^{n-1}}{n(n+2)}
=\frac{2^{n+1}}{n+2}-\frac{2^{n-1}}{n}\).

当 \( n \) 为偶数时，
\(c_n=\frac{a_{n-1}}{b_{n+1}} =\frac{n-1}{2^n}\).

对任意的正整数 \( n \) ，有
\(\sum_{k=1}^n c_{2k-1}
=\sum_{k=1}^n\left( \frac{2^{2k}}{2k+1}-\frac{2^{2k-2}}{2k-1} \right)
=\sum_{k=1}^n\left( \frac{4^k}{2k+1}-\frac{4^{k-1}}{2k-1} \right)
=\frac{4^n}{2n+1}-1\).

又
\(\sum_{k=1}^n c_{2k}
=\sum_{k=1}^n\frac{2k-1}{4^k}
=\frac14+\frac3{4^2}+\frac5{4^3}+\cdots+\frac{2n-1}{4^n}\).

记
\(A=\sum_{k=1}^n c_{2k}\)，
则
\(4A=1+\frac34+\frac5{4^2}+\cdots+\frac{2n-1}{4^{n-1}}\).

两式相减，得
\(3A=1+\frac2{4}+\frac2{4^2}+\cdots+\frac2{4^{n-1}}-\frac{2n-1}{4^n}\).

所以
\(3A=1+2\sum_{k=1}^{n-1}\frac1{4^k}-\frac{2n-1}{4^n}
=1+\frac23\left( 1-\frac1{4^{n-1}} \right)-\frac{2n-1}{4^n}
=\frac53-\frac{6n+5}{3\cdot4^n}\).
从而
\(\sum_{k=1}^n c_{2k} =A =\frac59-\frac{6n+5}{9\cdot4^n}\).

因此，
\(\sum_{k=1}^{2n}c_k
=\sum_{k=1}^n c_{2k-1}+\sum_{k=1}^n c_{2k}
=\frac{4^n}{2n+1}-1+\frac59-\frac{6n+5}{9\cdot4^n}
=\frac{4^n}{2n+1}-\frac49-\frac{6n+5}{9\cdot4^n}\).
\end{solution}
\begin{problem}
已知函数 \(f(x)=x^3+k\ln x\)（\(k\in\R\)），\(f'(x)\) 为 \(f(x)\) 的导函数.

\begin{enumerate}
\item 当 \(k=6\) 时，
\begin{enumerate}
\item 求曲线 \(y=f(x)\) 在点 \((1,f(1))\) 处的切线方程；
\item 求函数 \(g(x)=f(x)-f'(x)+\dfrac9x\) 的单调区间和极值；
\end{enumerate}
\item 当 \(k\ge-3\) 时，求证：对任意的 \(x_1,x_2\in[1,+\infty)\)，且 \(x_1>x_2\)，有 \(\dfrac{f'(x_1)+f'(x_2)}{2}>\dfrac{f(x_1)-f(x_2)}{x_1-x_2}\).
\end{enumerate}

\end{problem}

\begin{solution}
（1）[(1)]

（1）[(i)] 当 \( k=6 \) 时，
\(f(x)=x^3+6\ln x\)，\(f'(x)=3x^2+\frac6x\).
可得 \( f(1)=1 \) ，\( f'(1)=9 \) ，

所以曲线 \( y=f(x) \) 在点 \( (1,f(1)) \) 处的切线方程为
\(y-1=9(x-1)\)，
即
\(y=9x-8\).

（2）[(ii)] 依题意，
\(g(x)=x^3-3x^2+6\ln x+\frac3x\)，\(x\in(0,+\infty)\).

从而
\(g'(x)=3x^2-6x+\frac6x-\frac3{x^2} =\frac{3(x-1)^3(x+1)}{x^2}\).

令 \( g'(x)=0 \) ，解得 \( x=1 \) .

当 \( x \) 变化时， \( g'(x) \) ，\( g(x) \) 的变化情况如下表：

\examdisplayarray{}{c|ccc}{
x & (0,1) & 1 & (1,+\infty) \\
\hline
g'(x) & - & 0 & + \\
\hline
g(x) & \text{单调递减} & \text{极小值} & \text{单调递增}
}{}

所以，函数 \( g(x) \) 的单调递减区间为 \( (0,1) \) ，单调递增区间为 \( (1,+\infty) \) ；

\(g(1)=1-3+0+3=1\)，
故 \( g(x) \) 的极小值为 \( g(1)=1 \) ，无极大值.

（2）[(2)] 证明：由 \( f(x)=x^3+k\ln x \) ，得
\(f'(x)=3x^2+\frac{k}{x}\).

对任意的 \( x_1,x_2\in[1,+\infty) \) ，且 \( x_1>x_2 \) ，令
\(t=\frac{x_1}{x_2}\)（\(t>1\)），
则
\examdisplayaligned{
&\left( x_1-x_2 \right)\left( f'(x_1)+f'(x_2) \right)-2\left( f(x_1)-f(x_2) \right) \\
=&\left( x_1-x_2 \right)\left( 3x_1^2+\frac{k}{x_1}+3x_2^2+\frac{k}{x_2} \right)-2\left( x_1^3-x_2^3+k\ln\frac{x_1}{x_2} \right) \\
=&x_2^3\left( t^3-3t^2+3t-1 \right)+k\left( t-\frac1t-2\ln t \right).
}

令
\(h(x)=x-\frac1x-2\ln x\)，\(x\in[1,+\infty)\)，
则
\(h'(x)=1+\frac1{x^2}-\frac2x=\left( 1-\frac1x \right)^2>0\)（\(x>1\)）.

由此可得 \( h(x) \) 在 \( [1,+\infty) \) 上单调递增，所以当 \( t>1 \) 时，
\(t-\frac1t-2\ln t=h(t)>h(1)=0\).

因为 \( x_2\ge1 \) ，\( k\ge-3 \) ，且
\(t^3-3t^2+3t-1=\left( t-1 \right)^3>0\)，
所以
\examdisplayaligned{
&x_2^3\left( t^3-3t^2+3t-1 \right)+k\left( t-\frac1t-2\ln t \right) \\
\ge&\left( t^3-3t^2+3t-1 \right)-3\left( t-\frac1t-2\ln t \right) \\
=&t^3-3t^2+6\ln t+\frac3t-1.
}

由（1）（2）可知，当 \( t>1 \) 时，
\(g(t)=t^3-3t^2+6\ln t+\frac3t>g(1)=1\)，
即
\(t^3-3t^2+6\ln t+\frac3t-1>0\).

因此，
\(\left( x_1-x_2 \right)\left( f'(x_1)+f'(x_2) \right)-2\left( f(x_1)-f(x_2) \right)>0\).

所以，当 \( k\ge-3 \) 时，任意的 \( x_1,x_2\in[1,+\infty) \) ，且 \( x_1>x_2 \) ，有
\(\frac{f'(x_1)+f'(x_2)}{2}>\frac{f(x_1)-f(x_2)}{x_1-x_2}\).
\end{solution}
